\documentclass[12pt]{article}
\usepackage[utf8]{inputenc}
\usepackage[T1]{fontenc}
\usepackage{hyperref}
\title{Yamphe Language Resources for K-12 Curriculum Development}
\author{Compiled by Claude Sonnet 5.5 (claude-sonnet-5-5) for Andr\'as Kornai}
\date{}
\begin{document}
\maketitle

\section{Introduction}
Yamphe (ISO 639-3 yma; Glottocode yamp1243) is a problematic entry. The prompt lists its family as ``Bookkeeping'', which is not a linguistic family. Searches suggest that Yamphe is related to, or an alternative name for, Yamphu, a Kiranti (Sino-Tibetan) language of eastern Nepal. The Endangered Languages Project lists Yamphe as an alternative name for Yamphu. Wikipedia says the Tomyang (Chongka) dialect and Yamphe are both distinct from Yamphu, so the relationship is unresolved. I could not settle this, and the resources below are for Yamphu and may not apply to Yamphe.

\section{Primary Educational Resources}
I found no primers or school materials. Sources say Yamphu has no written tradition.

\section{Linguistic and Reference Materials}
\begin{itemize}
\item Rutgers, Roland (1998), \emph{Yamphu: grammar, texts and lexicon}, which describes the Hedangna dialect. I saw only a bookseller listing. \url{https://www.boekwinkeltjes.nl/b/231029811/}
\item The English Wikipedia article on Yamphu gives about 9,200 speakers (2011 census). \url{https://en.wikipedia.org/wiki/Yamphu_language}
\item The Endangered Languages Project entry for Yamphu lists Yamphe as an alternative name and classes the language as threatened. \url{https://endangeredlanguages.com/lang/3216}
\end{itemize}

\section{Religious and Community Resources}
I found none.

\section{Digital and Audio Resources}
I found none. A STEDT etymological listing is available. \url{https://stedt.berkeley.edu/~stedt-cgi/rootcanal.pl/group/31/1873}

\section{Curriculum Development Context}
Little exists. Before any curriculum work, the identity of Yamphe needs to be settled, since the Glottolog family field appears to be a data error. Any curriculum would be built with Kiranti and Nepali community bodies, none of which I found in these searches.
\end{document}